% What a word set again is met by at each of its ends: the boundary takes
% part only where it took part the first time.  A word whose boundary was
% cancelled at one end begins or ends in a plain character rather than in a
% ligature the boundary helped make, and is set again without one there.
\catcode`\{=1 \catcode`\}=2
\tracingonline=1 \showboxbreadth=99 \showboxdepth=99
\font\rip=trip \rip \hyphenchar\rip=`\- \lccode`1=`1
\hsize=200pt \parindent=0pt \pretolerance=-1 \tolerance=10000
\hyphenation{1-1}
\message{[A a word that ends the run]}
\setbox1=\vbox{\noindent\ 11 \par}\showbox1
\message{[B one whose boundary is cancelled]}
\setbox1=\vbox{\noindent\ 11\noboundary\ \par}\showbox1
\message{[C one begun after \string\noboundary]}
\setbox1=\vbox{\noindent\ \noboundary11 \par}\showbox1
\message{[done]}
\end
